% ============================================================
% SECTION 6 — INTEGRATION WITH THE BX3 FRAMEWORK
% ============================================================
\section{Integration with the BX3 Framework}
\label{sec:integration}

Recursive Spawning is not a standalone construct; it is embedded within the
three-layer accountability architecture defined by the BX3 Framework. Each
spawning event is choreographed across the Purpose Layer, Bounds Engine,
and Fact Layer, ensuring that agent lineage remains coherent, constrained,
and auditable from instantiation through termination.

% --- 6.1 Purpose Layer as Accountability Anchor ---
\subsection{Purpose Layer as Accountability Anchor}
\label{sec:integration:purpose}

The Purpose Layer serves as the \textit{accountability anchor} for every
spawned agent. Prior to any spawning event, the initiating agent must
declare the child agent's intent, scope, and success criteria encoded
as a Purpose Layer record. This record is immutable once written and
propagates down the chain such that every descendant inherits a
purpose-subgraph rooted in the original declaration.

Formally, let $p_\text{parent}$ be the purpose record of the spawning agent.
The spawning event produces a purpose record $p_\text{child}$ such that:

\begin{equation}
  p_\text{child} = \text{SpawnPurpose}(p_\text{parent}, \delta_\text{scope},
  \kappa_\text{criteria})
\end{equation}

where $\delta_\text{scope}$ denotes the intentional deviation (narrowing,
lateral shift, or extension) from the parent's scope, and
$\kappa_\text{criteria}$ defines the measurable termination conditions.
The Purpose Layer stores this record as an append-only entry in the
agent's purpose ledger. No agent may spawn without a Purpose Layer record;
the Bounds Engine enforces this prerequisite before granting a spawn
permit.

The anchor property is transitive: because every agent in the chain holds
a purpose record that traces, via parent pointers, to a Purpose Layer
entry at depth zero, the entire lineage is anchored to an authoritative
intent declaration. This prevents purpose drift by making the original
intent permanently retrievable and binding on all descendants.

% --- 6.2 Bounds Engine Depth Enforcement ---
\subsection{Bounds Engine Depth Enforcement}
\label{sec:integration:bounds}

The Bounds Engine provides the \textit{operational constraints} that keep
Recursive Spawning within defined limits. At spawn time, the Bounds Engine
evaluates three constraint categories:

\begin{enumerate}
  \item \textbf{Depth limit}: The cumulative spawn depth $d_\text{cum}$
    must satisfy $d_\text{cum} \leq D_\text{max}$, where $D_\text{max}$
    is a workspace-global constant enforced by the Bounds Engine policy.
    Violations result in immediate rejection of the spawning event.
  \item \textbf{Scope aperture}: The child's scope aperture
    $\alpha_\text{child}$ must be a proper subset of $\alpha_\text{parent}$
    or a laterally adjacent scope within the same capability manifold.
    Lateral adjacency is defined by the Bounds Engine's scope adjacency
    matrix $\mathbf{M}_\text{adj}$.
  \item \textbf{Resource budget}: Each spawn consumes a configurable
    resource debit $r_\text{spawn}$ from the parent's allocated budget
    $B_\text{parent}$. Spawning proceeds only if
    $B_\text{parent} - r_\text{spawn} \geq 0$.
\end{enumerate}

These constraints are evaluated atomically at the moment of the spawning
event. The Bounds Engine issues a signed constraint attestation — a
cryptographic manifest recording the constraint values at spawn time —
which is stored alongside the Fact Layer record for the event. This
ensures that post-hoc forensic analysis can determine the exact constraint
state that governed any given spawn.

% --- 6.3 Fact Layer Forensic Ledger ---
\subsection{Fact Layer Forensic Ledger}
\label{sec:integration:fact}

The Fact Layer functions as the \textit{forensic ledger} for the entire
spawning chain. Every spawning event, state transition, and termination
is recorded as an append-only, tamper-evident fact entry. The Fact Layer
entry for a spawn event contains:

\begin{itemize}
  \itemsep0em
  \itemsep0em
  \itemsep0em
  \item Spawn timestamp $t_\text{spawn}$ (monotonically increasing)
  \item Parent agent identifier $\text{ID}_\text{parent}$
  \item Child agent identifier $\text{ID}_\text{child}$ (assigned by the
    Fact Layer registrar)
  \item Purpose Layer record hash $h(p_\text{child})$
  \item Bounds Engine constraint attestation hash $h(\text{attest})$
  \item Chain depth $d_\text{child} = d_\text{parent} + 1$
  \item Immutable parent pointer $P_\text{parent}$ (hash-chained)
\end{itemize}

The parent pointer $P_\text{parent}$ is itself a Fact Layer entry hash,
creating a cryptographically chain-linked sequence from any agent back
to its chain origin. This structure satisfies the requirements of a
blockchain-style ledger without requiring distributed consensus: the
Fact Layer is a centralized append-only log with hash-chain integrity
enforced by the BX3 runtime. Any attempt to alter a historical entry
breaks the hash chain at that point, making tampering detectable by
standard chain-integrity verification.

Together, the three layers produce a spawning event in which:
purpose declares intent, the Bounds Engine permits or denies based on
operational constraints, and the Fact Layer records the immutable
deterministic record of what was permitted and why. No layer can
subvert another without breaking its own integrity guarantees.


% ============================================================
% SECTION 7 — CORRECTNESS PROPERTIES
% ============================================================
\section{Correctness Properties}
\label{sec:correctness}

Recursive Spawning satisfies three correctness properties of primary
importance to agentic system design: (i) purpose drift is prevented,
(ii) chain integrity is maintained across the full lineage, and
(iii) spawning terminates at a defined depth boundary. We formalize
each property and outline the proof structure.

% --- 7.1 Drift Prevention ---
\subsection{Drift Prevention}
\label{sec:correctness:drift}

\textit{Purpose drift} is the failure mode in which an agent's behavior
deviates from its declared intent over successive spawn generations,
such that descendants no longer serve the original purpose. Drift can
arise through scope creep, unbounded lateral expansion, or the gradual
loss of intent context across spawn links.

\begin{theorem}[Drift Prevention]
  \label{th:drift}
  In a Recursive Spawning system governed by the BX3 Framework, no
  agent in any spawn chain can exhibit purpose drift relative to its
  chain-origin purpose record.
\end{theorem}

\begin-proof}
  We prove Theorem~\ref{th:drift} by demonstrating that the conjunction
  of three conditions — each guaranteed by a distinct BX3 layer —
  makes drift structurally impossible.

  \begin{condition}
    \label{cond:immutable}
    \textbf{Immutable parent pointer.} The parent pointer $P_\text{parent}$
    stored in the Fact Layer for any agent is a hash-chained reference
    to the parent's Fact Layer entry and cannot be altered retroactively
    without breaking the chain-integrity hash. Therefore, the parentage
    of any agent is permanently and verifiably recorded.
  \end{condition}

  \begin{condition}
    \label{cond:forensic}
    \textbf{Forensic ledger.} The Fact Layer append-only log records
    the full spawn event, purpose record, and constraint attestation
    for every generation. Any deviation from the prescribed purpose is
    detectable by cross-referencing the agent's executed scope against
    its declared $p_\text{child}$ scope aperture.
  \end{condition}

  \begin{condition}
    \label{cond:terminus}
    \textbf{Chain terminates at Purpose Layer depth zero.} The spawn chain
    is anchored at depth zero to a Purpose Layer record with no parent.
    All descendant purpose records are intentionally derived from this
    anchor via the SpawnPurpose relation. Because SpawnPurpose always
    produces a $p_\text{child}$ that is a scoped variant of its parent,
    and because the Bounds Engine enforces scope-aperture proper-subset
    or lateral-adjacency constraints, the descendant's scope cannot
    drift to an unrelated domain.
  \end{condition}

  \begin{claim}
    \label{clm:drift-impossible}
    Drift is impossible under conditions~\ref{cond:immutable},
    \ref{cond:forensic}, and~\ref{cond:terminus}.
  \end{claim}

  \begin{subproof}
    Suppose for contradiction that an agent $A_k$ at depth $k$ exhibits
    drift: its executed scope $\sigma_k$ is not a scoped variant of the
    anchor purpose $p_0$. By Condition~\ref{cond:terminus}, $A_k$'s
    purpose record $p_k$ must be a scoped variant of $p_0$. By
    Condition~\ref{cond:forensic}, the Fact Layer records the spawn
    event linking $A_k$ to its parent $A_{k-1}$ and storing $h(p_k)$.
    The executed scope $\sigma_k$ can be cross-verified against $p_k$
    at any time using the forensic ledger. If $\sigma_k \not\subseteq
    \sigma(p_k)$, the drift is detected. Since the Fact Layer is
    append-only, $p_k$ cannot be retroactively rewritten to conceal
    the drift (Condition~\ref{cond:immutable}). Therefore no undetected
    drift can exist; any drift that does occur is immediately surfaced
    as aFact Layer anomaly. Drift is thus structurally prevented, not
    merely detected after the fact.
  \end{subproof}
\end-proof}

The key structural insight is that drift prevention does not rely on
runtime monitoring alone. It is guaranteed by the immutable
parent-pointer chain (Condition~\ref{cond:immutable}), the forensic
ledger that makes any divergence immediately observable
(Condition~\ref{cond:forensic}), and the Purpose Layer anchoring that
constrains the scope-manifold to intentional descendants only
(Condition~\ref{cond:terminus}).

% --- 7.2 Chain Integrity ---
\subsection{Chain Integrity}
\label{sec:correctness:integrity}

\begin{theorem}[Chain Integrity]
  \label{th:integrity}
  For any agent $A_k$ in a Recursive Spawning chain, the full lineage
  $\{A_0, A_1, \dots, A_k\}$ is retrievable, unmodified, and
  cryptographically verifiable via Fact Layer hash-chain traversal.
\end{theorem}

\begin-proof}
  By construction, each Fact Layer entry $e_i$ for $i \geq 1$ contains
  $P_\text{parent} = h(e_{i-1})$. The chain traversal algorithm
  $\text{VerifyChain}(A_k)$ proceeds as follows:

  \begin{enumerate}
    \item Obtain $e_k$ from the Fact Layer registry using
      $\text{ID}_{A_k}$.
    \item Compute $\hat{h} = h(e_{k-1})$ from the stored parent pointer
      $P_\text{parent}$ in $e_k$.
    \item Compare $\hat{h}$ against the stored hash of $e_{k-1}$. If
      they differ, the chain is broken and verification fails.
    \item Recurse to $e_{k-1}$ and repeat until $e_0$ is reached.
  \end{enumerate}

  Because each entry is append-only and the hash function is
  collision-resistant, any modification to a historical entry produces
  a hash mismatch at the point of modification and all subsequent
  entries. Therefore, successful chain traversal from $A_k$ to $A_0$
  is both necessary and sufficient for chain integrity. The chain
  terminates at $A_0$, which is identified by the special marker
  $P_\text{parent} = \varnothing$, indicating a Purpose Layer
  depth-zero origin.
\end-proof}

% --- 7.3 Termination ---
\subsection{Termination}
\label{sec:correctness:termination}

\begin{theorem}[Termination]
  \label{th:termination}
  Every Recursive Spawning chain terminates: there exists a finite
  maximum depth $D_\text{max}$ such that no agent may spawn at depth
  $d > D_\text{max}$. All spawned agents either terminate autonomously
  upon satisfaction of their $\kappa_\text{criteria}$ or are terminated
  by the Bounds Engine upon reaching $D_\text{max}$.
\end{theorem}

\begin-proof}
  The Bounds Engine enforces the depth constraint
  $d_\text{cum} \leq D_\text{max}$ as a mandatory precondition for every
  spawning event (Section~\ref{sec:integration:bounds}). A spawn request
  for which $d_\text{child} > D_\text{max}$ is rejected atomically
  before the Fact Layer entry is created. Consequently, no agent can
  exist at a depth exceeding $D_\text{max}$.

  For any agent $A_i$ at depth $d_i \leq D_\text{max}$, the agent's
  purpose record $p_i$ encodes termination criteria $\kappa_i$. The
  Fact Layer records whether $A_i$ reached its $\kappa_i$ milestone or
  was Bounds-Engine-terminated due to depth exhaustion. In either case,
  the agent ceases to execute and its spawn budget is zeroed. Since the
  depth counter is strictly incremented at each spawn and bounded above
  by $D_\text{max}$, and since each agent terminates in finite time
  by the BX3 runtime model, the total number of agents in any chain is
  bounded by $D_\text{max} + 1$, a finite constant. Hence, all chains
  are finite and termination is guaranteed.
\end-proof}

These three properties — drift prevention, chain integrity, and
termination — form the correctness triad of Recursive Spawning under
the BX3 Framework. Together they ensure that agentic systems built on
this substrate can spawn complex, multi-generational workflows with
formal guarantees of intent preservation, auditable lineage, and
bounded resource consumption.
